% Generated draft: requires mathematical and source review.
\documentclass[11pt]{article}
\usepackage[margin=25mm]{geometry}
\usepackage{fontspec}
\setmainfont{Latin Modern Roman}
\usepackage{amsmath,amssymb,mathtools,unicode-math}
\setmathfont{Latin Modern Math}
\usepackage{enumitem,needspace,xurl,hyperref}
\setlength{\parindent}{0pt}
\setlength{\parskip}{5pt}
\setlength{\emergencystretch}{3em}
\newcommand{\N}{\mathbb{N}}
\newcommand{\Z}{\mathbb{Z}}
\newcommand{\Q}{\mathbb{Q}}
\newcommand{\R}{\mathbb{R}}
\newcommand{\C}{\mathbb{C}}
\newcommand{\F}{\mathbb{F}}
\newcommand{\A}{\mathbb{A}}
\newcommand{\PP}{\mathbb{P}}
\newcommand{\E}{\mathbb{E}}
\newcommand{\eps}{\varepsilon}
\newcommand{\dd}{\,\mathrm{d}}
\begin{document}
\section*{UnsolvedMath ID: 30000922}
\textbf{Source problem code:} catalogue entry
\subsection{Definitions and mathematical statement}
Let \(D=\{z\in\C:|z|<1\}\), with boundary equipped with arc-length measure. An inner function \(B\) is a bounded holomorphic function on \(D\) whose radial boundary values have modulus one almost everywhere. A Blaschke product \(I\) is interpolating if its zero sequence \((a_n)\) is an interpolating sequence for \(H^\infty(D)\): every bounded scalar sequence \((c_n)\) is of the form \(c_n=h(a_n)\) for some bounded holomorphic \(h\).

For a real-valued \(v\in L^\infty(\partial D)\), let \(\widetilde v\) denote the boundary function of the harmonic conjugate of its Poisson extension, conventionally normalized to vanish at zero. The record asks the following two questions for every inner function \(B\):
\begin{enumerate}
\item Does there exist an interpolating Blaschke product \(I\) and real-valued \(v\in L^\infty(\partial D)\) such that
\[
\operatorname{Arg}B-\operatorname{Arg}I=\widetilde v\quad\text{almost everywhere on }\partial D?
\]
\item Does there exist an interpolating Blaschke product \(I\) and real-valued \(u,v\in L^\infty(\partial D)\) such that
\[
\operatorname{Arg}B-\operatorname{Arg}I=u+\widetilde v\quad\text{almost everywhere},
\qquad \|u\|_\infty<\pi/2?
\]
\end{enumerate}
The choices of \(I\) and the boundary functions may depend on \(B\) and may differ between the two questions. The second question retains the strict bound \(\pi/2\), rather than an arbitrary bounded error.

The source does not fix boundary-argument branches, equality of real-valued lifts versus equality modulo \(2\pi\), harmonic-conjugate additive constants, or constant unimodular factors in \(I\). These conventions must be checked; no principal-argument requirement is silently imposed. The normalized conjugate above is an explicit convention for interpreting the symbol, not a source-verified normalization.
\subsection{Short English statement}
Can every inner-function boundary argument be matched by an interpolating Blaschke-product argument up to a conjugate of a bounded real function, either exactly or with an additional bounded error of norm strictly below \(\pi/2\)?
\subsection{Sources}
\begin{itemize}
\item Local UnsolvedMath registry, permanent ID \(30000922\), source code catalogue entry, ``Blaschke Approximation of Inner-Function Arguments.'' No individual proposer is supplied.
\item \emph{Mini-Workshop: Complex Approximation and Universality} (2008), Oberwolfach Reports, DOI 10.4171/owr/2008/06, pp. 297--346: \url{https://doi.org/10.4171/owr/2008/06}.
\item Supplied reference, ``argument approximation paper'': \url{https://msp.org/pjm/2006/228-1/pjm-v228-n1-p05-p.pdf}.
\item Supplied literature triage, checked 2026-08-21: status ``open.'' It says the report and a Pacific Journal paper present these questions as open and reports no located proof or counterexample. It calls that paper ``later,'' although the supplied URL contains 2006 while the report is dated 2008; this chronology is unresolved. This dated assessment derives from supplied OWR curation; time-sensitive status and cited claims require independent verification.
\end{itemize}
\end{document}
